\chapter{User Stories}
\ChapterLabel{UserStories}
\index{user story}
\index{acceptance criterion}
\section{Chapter Responsibility}
\label{Section:UserStoriesResponsibility}
Capture who needs a result and why that result has value.

\noindent\textbf{Expected result:} A reviewable statement of stakeholder value and acceptance intent.

\section{Revision 0.2 Scope Notice}
\label{Section:DeferredUserStories}
\begin{fsaNote}[Project drafts and retained template reference]
The two project stories below are Draft for review. Joel Murphy (co-owner from Version 0.2: Jason Lefler) approved
their insertion on November 21, 2026; that approval does not baseline
requirements, establish implementation, or approve validation results.
Section 4.5 uses the template register format with the two project story
rows replacing the original examples. Legacy template IDs remain only as
reference targets for unchanged example chapters; they do not describe
project behavior. Later requirements, architecture,
tests, verification, validation, and release content remain planned or pending
under \hyperref[Section:ProjectScope]{the revision scope}.
\end{fsaNote}

\section{How to Use This Chapter}
Capture stakeholder goals in accessible language. Stories express value and
acceptance intent; precise obligations belong in Requirements. Give each story
a stable ID, a priority, and links to one or more requirements that implement it.

\section{Project User Stories --- Drafts for Review}
\label{Section:ProjectDraftStories}
The two project drafts are defined in
\hyperref[Table:UserStories]{the User Story Register} below, using the same
ID, Story and Acceptance Criteria, Priority, and Requirement(s) columns as
the original template. Both retain revision 0.2 and Draft for review status.
FTD2XX is the proposed library for FTDI's D2XX Classic API. No library
version, platform compatibility, or implementation result is implied. The
D2XX Programmer's Guide v1.3 is the approved documentation baseline; actual
driver/library version remains TBD (decision co-owners: Joel Murphy and Jason Lefler).
These stories follow \hyperref[Skill:AIS-05]{AIS-05} and its
\hyperref[Section:AIUserStoryContext]{User-Story Context}.

\paragraph{Related draft use cases.}
\UserStoryRef{US-001} maps to \UseCaseRef{UC-001};
\UserStoryRef{US-002} maps to \UseCaseRef{UC-002}.
See \hyperref[Figure:UseCaseDiagram]{the linked UML diagram and use-case tables}.
Requirement proposals and planned tests are linked in \hyperref[Section:ProjectRequirements]{Project Requirements}; review decisions remain with Joel Murphy and Jason Lefler.

\section{User Story Register}
\label{Section:ProjectUserStoryRegister}
The template example rows have been replaced by
\UserStoryRef{US-001} and \UserStoryRef{US-002}. All unresolved decisions in
each row, including priority and later lifecycle allocation, are assigned to
Joel Murphy and Jason Lefler. Candidate requirement allocations are linked below; validation
execution and the remaining lifecycle work remain planned or pending.

\begingroup
\small
\setlength{\emergencystretch}{2em}
\begin{tabularx}{\linewidth}{L{1.9cm}YL{1.8cm}L{3.4cm}}
\caption{User Story Register --- Project Drafts\label{Table:UserStories}}\\
\toprule
\textbf{ID} & \textbf{Story and Acceptance Criteria} & \textbf{Priority} & \textbf{Requirement(s)}\\
\midrule
\endfirsthead
\multicolumn{4}{l}{\small\itshape User Story Register --- Project Drafts (continued)}\\
\toprule
\textbf{ID} & \textbf{Story and Acceptance Criteria} & \textbf{Priority} & \textbf{Requirement(s)}\\
\midrule
\endhead
\midrule
\multicolumn{4}{r}{\small\itshape Continued on next page}\\
\endfoot
\bottomrule
\endlastfoot
\UserStoryID{US-001}\phantomsection\label{Section:AcquireSensorDataStory}
&\par As an Embedded Engineer, I need to receive the data
supplied through the FTDI245R interface and make it available to the
downstream processing consumer, so that I can process the acquired
information on the host.
Acceptance: A known source byte sequence reaches the consumer
unchanged and in order under agreed conditions.
& Must
& \RequirementRef{REQ-F-008}; \RequirementRef{REQ-F-009}; \RequirementRef{REQ-F-010}; \RequirementRef{REQ-I-002}; \RequirementRef{REQ-D-002};
\cite{FTDI_TN167_v1_0}, \cite{FTDI_D2XX_v1_3}. \\
\UserStoryID{US-002}\phantomsection\label{Section:AcquisitionInterruptionsStory}
&\par As an Embedded Engineer, I need the pipeline to distinguish
completed reads, partial or empty timeout results, and read failures, so that
interruptions do not crash or corrupt the pipeline or cause invalid data to
reach downstream processing.
Acceptance: Completed reads are recognized using status and
returned count.
& Should
& \RequirementRef{REQ-F-004}; \RequirementRef{REQ-F-005}; \RequirementRef{REQ-F-006}; \RequirementRef{REQ-F-007}; \RequirementRef{REQ-F-011}; \RequirementRef{REQ-F-012}; \RequirementRef{REQ-F-013}; \RequirementRef{REQ-Q-002}; \RequirementRef{REQ-D-002} \cite{FTDI_D2XX_v1_3}.
\end{tabularx}
\endgroup

\paragraph{Legacy template cross-link targets.}
The following IDs identify superseded template examples, not active project
stories: \UserStoryID{TEMPLATE-US-001} (submit a service request),
\UserStoryID{TEMPLATE-US-002} (view request status), and
\UserStoryID{US-003} (export request history). These reference targets are
retained solely for links from the unchanged template example chapters;
their example requirements and validation records are not allocated to
\UserStoryRef{US-001} or \UserStoryRef{US-002}.

\paragraph{Draft traceability status.}
\hyperref[Section:DraftStoryTraceability]{The planned project traceability register}
provides reciprocal links to both project stories. Requirements, use cases,
architecture, tests, verification, validation execution, and release remain
planned or pending, with unresolved decisions assigned to Joel Murphy and Jason Lefler.

\section{Review Rules}
Confirm that every story identifies a beneficiary, outcome, acceptance
criterion, priority, and at least one linked requirement. Maintain reciprocal
story links in the traceability register.

The project drafts retain explicit pending requirement allocations until
requirements are developed and reviewed. Their presence does not establish
complete baseline traceability; the open decisions are assigned to Joel Murphy and Jason Lefler.
